\chapter{多项式}\label{ch:b1:poly}

\omterm{def:b1:poly:def}{多项式}是代数学家最喜爱的函数——只是
这里并不把它们当作函数，而是把它们看作
不定元 $X$ 上的形式表达式，按
交换环的规则做加法与乘法。这一理论与
\cref{ch:b1:arith}惊人地平行：欧几里得除法、\omterm{thm:b1:arith:gcd}{最大公因数}与贝祖关系、
不可约元与唯一分解。全文中，$K$
表示 $\Q$、$\R$ 或 $\C$。

\section{环 $K[X]$}

\begin{definition}[多项式，次数]\label{def:b1:poly:def}
称系数在 $K$ 中的\emph{多项式}\index{多项式}为
形式和：
\[
P = a_0 + a_1 X + a_2 X^2 + \dots + a_n X^n
= \sum_{k} a_k X^k,
\]
其中 $a_k \in K$ 从某一指标起全为零。按自然的
加法及下式定义的乘法，
\[
\Bigl(\sum_i a_i X^i\Bigr)\Bigl(\sum_j b_j X^j\Bigr)
= \sum_k \Bigl(\sum_{i+j=k} a_i b_j\Bigr) X^k,
\]
\omterm{def:b1:logic:sets}{集合} $K[X]$ 是交换环。\emph{次数}$\deg P$
（$P \neq 0$）是最大的 $n$，使 $a_n \neq 0$；$a_n$ 称为
\emph{首项系数}（$P$ 为\emph{首一}\index{首一
多项式}，当 $a_n = 1$ 时），并约定 $\deg 0 = -\infty$。
每个多项式都经代入确定一个函数 $x \mapsto P(x)$，
它定义在 $K$ 上。
\end{definition}

\begin{proposition}[次数的运算法则；整环]\label{prop:b1:poly:degree}
设 $P, Q \in K[X]$，有：
\[
\deg(P + Q) \leq \max(\deg P, \deg Q),
\qquad
\deg(PQ) = \deg P + \deg Q .
\]
由此 $K[X]$ 是\omterm{def:b1:structures:field}{整环}，且其单位恰为
非零常数。
\end{proposition}

\begin{proof}
和的运算法则显然成立（超出最大次数的系数均为零）。对于
乘积，设 $a_m$ 与 $b_n$ 为首项系数，则
$X^{m+n}$ 在 $PQ$ 中的系数为 $a_m b_n \neq 0$（$K$ 是
域，故为\omterm{def:b1:structures:field}{整环}），且更高次的系数全为零。
若 $P, Q \neq 0$，则 $\deg PQ = \deg P + \deg Q \geq 0$，故 $PQ \neq
0$，即为\omterm{def:b1:structures:field}{整环}。若 $PQ = 1$，则 $\deg P + \deg Q = 0$ 迫使
$\deg P = \deg Q = 0$。于是可逆元恰为可逆常数，
即 $K^*$ 中的全部元素。
\end{proof}

\begin{theorem}[欧几里得除法]\label{thm:b1:poly:division}
设 $A, B \in K[X]$，$B \neq 0$。则恰有一对\omterm{def:b1:poly:def}{多项式} $(Q,
R)$ 满足
\[
A = BQ + R, \qquad \deg R < \deg B .
\]
\end{theorem}

\begin{proof}
\emph{存在性}，对 $\deg A$ 作强归纳法。若 $\deg A < \deg
B$，取 $(Q, R) = (0, A)$。否则写 $A = a X^m + \dots$，$B =
b X^n + \dots$ 且 $m \geq n$；\omterm{def:b1:poly:def}{多项式} $A_1 = A - \frac ab
X^{m-n} B$ 的次数为 $< m$（首项相消），故由
归纳假设 $A_1 = BQ_1 + R$，其中 $\deg R < \deg B$，且 $A = B(Q_1 +
\frac ab X^{m-n}) + R$。

\emph{唯一性}：若 $BQ + R = BQ' + R'$，则 $B(Q - Q') = R' - R$
且 $\deg(R' - R) < \deg B$。由次数运算法则可知 $Q - Q'
= 0$，于是 $R = R'$。
\end{proof}

\begin{example}\label{ex:b1:poly:division}
将 $A = X^4 + X^3 - 2X + 1$ 除以 $B = X^2 + 1$：
\[
X^4 + X^3 - 2X + 1 = (X^2 + 1)(X^2 + X - 1) + (-3X + 2).
\]
（计算：先减去 $X^2 B$，再减去 $X B$，最后减去 $-B$；余式
$-3X + 2$ 的次数为 $1 < 2$。）
\end{example}

\begin{method}[霍纳法]\label{met:b1:poly:horner}
欲求 $P = a_nX^n + \dots + a_0$ 在 $x$ 处的值，或将 $P$ 除以 $X - x$，
则应避免计算幂：从左到右依次读取系数，
并反复进行\emph{乘以 $x$，再加上下一个
系数}：
\[
b_n = a_n, \qquad b_{k} = a_{k} + x\,b_{k+1}
\quad (k = n-1, \dots, 0) .
\]
于是 $b_0 = P(x)$，而前面的 $b_k$ 即为商的
系数：$P = (X - x)(b_nX^{n-1} + \dots + b_1) + b_0$
（展开比较即可）。例：$P = X^4 - 5X^3 + 6X^2 + 4X - 8$
取 $x = 2$：$b$ 依次为 $1, -3, 0, 4, 0$，故 $P(2) = 0$，且 $P
= (X-2)(X^3 - 3X^2 + 4)$——一行即可，免去冗长的
除法，所需乘法次数仅为 $n$，而朴素求值则需 $\approx n^2/2$
次。在同一点上反复使用该格式，
可求出\omterm{def:b1:poly:derivative}{重数}（参见
\cref{ex:b1:poly:multexample}）。
\end{method}

\begin{remark}[{$K[X]$ 的算术}]
有了欧几里得除法之后，
\cref{ch:b1:arith}中的全部算术理论均可移植到 $K[X]$，证明相同，
以次数充当绝对值的角色：\omterm{thm:b1:arith:gcd}{最大公因数}（规范为\omterm{def:b1:poly:def}{首一}）、
扩展\omterm{met:b1:arith:euclid}{欧几里得算法}、贝祖恒等式、高斯引理、
不可约\omterm{def:b1:poly:def}{多项式}与唯一分解。这些移植而来的结论可以自由使用，
\cref{exo:b1:poly:6}对其中之一作了演练。
\end{remark}

\section{根}

\begin{theorem}[因式定理]\label{thm:b1:poly:factor}
设 $P \in K[X]$，$a \in K$。$P$ 除以
$X - a$ 所得余式为常数 $P(a)$。特别地
\[
P(a) = 0 \iff (X - a) \mid P .
\]
更一般地，互异的根 $a_1, \dots, a_r$ 给出 $P$ 的
分解式 $P = (X - a_1)\cdots(X - a_r)\, Q$。
\end{theorem}

\begin{proof}
作带余除法：$P = (X - a) Q + R$，$\deg R < 1$，故 $R$ 为常数
$c$；代入 $X = a$（代入保持和与积）
得 $P(a) = c$，等价性由此得出。对若干根，
对 $r$ 用归纳法，$r = 1$ 的情形即为刚证明的等价性。
假设结论对 $r - 1$ 个根成立，设 $a_1, \dots, a_r$
为 $P$ 的互异根。写 $P = (X - a_1)Q_1$；对每个 $i
\geq 2$，代入 $a_i$：
\[
0 = P(a_i) = (a_i - a_1)\,Q_1(a_i),
\qquad a_i - a_1 \neq 0 ,
\]
由于 $K$ 无零因子，$Q_1(a_i) = 0$，于是 $r - 1$
个互异的点 $a_2, \dots, a_r$ 都是 $Q_1$ 的根。
归纳假设给出 $Q_1 = (X - a_2)\cdots(X - a_r)\,Q$，
代回即得结论。
\end{proof}

\begin{corollary}[$n$ 次多项式至多有 $n$ 个根]\label{cor:b1:poly:nroots}
非零\omterm{def:b1:poly:def}{多项式} $P \in K[X]$ 次数为 $n$，至多有 $n$ 个
互异根（在 $K$ 中）。从而，次数 $\leq n$ 的\omterm{def:b1:poly:def}{多项式}在
$n + 1$ 个互异点处取零值，则为零\omterm{def:b1:poly:def}{多项式}，两个
次数 $\leq n$ 的\omterm{def:b1:poly:def}{多项式}在 $n+1$ 个点处取值相同，则二者相等。
\end{corollary}

\begin{proof}
若 $a_1, \dots, a_r$ 为互异根，由\cref{thm:b1:poly:factor}
可得 $P = (X-a_1)\cdots(X-a_r) Q$，故 $n = \deg P \geq r$。两个
推论分别用反证法和作差得到。
\end{proof}

\begin{example}[辅助多项式技巧]\label{ex:b1:poly:auxiliary}
设 $P$ 为次数 $\leq n$ 的\omterm{def:b1:poly:def}{多项式}，满足
\[
P(k) = \frac{k}{k+1} \qquad (k = 0, 1, \dots, n) ;
\]
由下文的\omterm{thm:b1:poly:lagrange}{拉格朗日插值}，该\omterm{def:b1:poly:def}{多项式}存在且唯一。求
$P(n+1)$？去分母：\omterm{def:b1:poly:def}{多项式} $Q = (X+1)P - X$ 的
次数 $\leq n + 1$，且在 $n + 1$ 个点 $0, 1,
\dots, n$ 处为零，由\cref{thm:b1:poly:factor}可知：
\[
Q = c\,X(X-1)(X-2)\cdots(X-n)
\]
某个常数 $c$。在 $Q$ 的值可独立算出的点处求值：
取 $X = -1$，$Q(-1) = 0 \cdot P(-1) + 1 = 1$，而乘积
为 $(-1)(-2)\cdots(-1-n) = (-1)^{n+1}(n+1)!$，故 $c =
\frac{(-1)^{n+1}}{(n+1)!}$。再取 $X = n + 1$ 求值：
\[
(n+2)\,P(n+1) - (n+1) = Q(n+1) = c\,(n+1)! = (-1)^{n+1} ,
\]
故 $P(n+1) = \dfrac{(n+1) + (-1)^{n+1}}{n+2}$，其值为 $1$
（$n$ 为奇数时），$\frac{n}{n+2}$（$n$ 为偶数时）——插值
\omterm{def:b1:poly:def}{多项式}\emph{并不}延续规律
$\frac{n+1}{n+2}$。值得记住的技巧：把已知数据编码为
某辅助\omterm{def:b1:poly:def}{多项式}的根，在已知数据之外的一点处确定未知常数，
再取值求得结论。
\end{example}

\begin{definition}[导数，重数]\label{def:b1:poly:derivative}
$P = \sum a_k X^k$ 的\emph{形式导数}为 $P' = \sum_{k
\geq 1} k\,a_k X^{k-1}$。它满足通常的法则 $(P+Q)' = P' +
Q'$，$(PQ)' = P'Q + PQ'$（在单项式上验证，
并由线性性推广得到）。设 $a$ 为 $P$ 的根，它的
\emph{重数}\index{重数}为 $m \geq 1$，当
$(X-a)^m \mid P$ 但 $(X-a)^{m+1} \nmid P$ 时；
$m = 1$ 时称根为\emph{单根}，$m \geq 2$ 时为\emph{重根}。
\end{definition}

\begin{proposition}[用导数刻画重数]\label{prop:b1:poly:multiplicity}
$a$ 是 $P$ 的根且\omterm{def:b1:poly:derivative}{重数} $\geq m$，当且仅当
\[
P(a) = P'(a) = \dots = P^{(m-1)}(a) = 0 .
\]
特别地，$a$ 为 $P$ 的重根当且仅当 $P(a) =
P'(a) = 0$.
\end{proposition}

\begin{proof}
写 $P = (X - a)^m Q + R$，其中 $R$ 是除以
$(X-a)^m$ 所得的余式，$\deg R < m$。对 $k \leq m - 1$ 次求导
并在 $a$ 处取值：第一项贡献 $0$（每次求导都保留因子
$(X-a)$），故 $P^{(k)}(a) = R^{(k)}(a)$。

现设\omterm{def:b1:poly:def}{多项式} $R$ 的次数 $< m$，它由 $R(a), R'(a),
\dots, R^{(m-1)}(a)$ 决定：写 $R = \sum_{k < m} c_k (X - a)^k$
（这可行：将 $X = (X - a) + a$ 代入并展开幂），
可得 $R^{(k)}(a) = k!\, c_k$。于是所有 $P^{(k)}(a) = 0$（$k < m$）
$\iff$ 所有 $c_k = 0$ $\iff$ $R = 0$ $\iff$ $(X-a)^m \mid P$。
\end{proof}

\begin{example}[计算重数]\label{ex:b1:poly:multexample}
求 $2$ 作为\omterm{def:b1:poly:def}{多项式} $P = X^4 - 5X^3 + 6X^2
+ 4X - 8$ 的根的\omterm{def:b1:poly:derivative}{重数}？逐阶计算导数在 $2$ 处的值：
\[
P(2) = 16 - 40 + 24 + 8 - 8 = 0, \qquad
P'(2) = 32 - 60 + 24 + 4 = 0,
\]
\[
P''(2) = 48 - 60 + 12 = 0, \qquad
P'''(2) = 48 - 30 = 18 \neq 0
\]
（其中 $P' = 4X^3 - 15X^2 + 12X + 4$，$P'' = 12X^2 - 30X + 12$，
$P''' = 24X - 30$）。前三个值为零，随后一个非零：
故\omterm{def:b1:poly:derivative}{重数}恰为 $3$。作除法，$P = (X - 2)^3(X + 1)$——
读者可通过展开 $(X-2)^3 = X^3 - 6X^2 + 12X - 8$ 并
再乘以 $X + 1$ 进行验证。要点在于：\omterm{def:b1:poly:derivative}{重数}可由
\emph{取值}直接读出，无需因式分解——这恰好是
在因式分解难以求得时检测\omterm{def:b1:poly:derivative}{重数}的方法。
\end{example}

\begin{example}[用最大公因数检测重根]\label{ex:b1:poly:gcdmultiple}
当未知任何根时，\cref{prop:b1:poly:multiplicity}仍
给出一个\emph{整体}的重根判别法：$a$ 为
$P$ 的重根，当且仅当它是 $P$ 与 $P'$ 的公共根。于是 $P$ 有
重根（在 $\C$ 中），当且仅当 $\gcd(P, P') \neq 1$——这可由
\omterm{met:b1:arith:euclid}{欧几里得算法}算出，无需求解任何方程。例：$P =
X^3 - 3X + 2$，$P' = 3X^2 - 3 = 3(X - 1)(X + 1)$。检验
$\pm1$（$P'$ 的根）代入 $P$：$P(1) = 0$，而 $P(-1) = 4$，故
\[
\gcd(P, P') = X - 1 :
\]
根 $1$ 为重根；连续除两次，得 $P = (X - 1)^2(X + 2)$。
\omterm{thm:b1:arith:gcd}{最大公因数}甚至给出全部重根的\omterm{def:b1:logic:sets}{集合}，每个重根的
\omterm{def:b1:poly:derivative}{重数}都降低一——这一事实被每个计算机代数
系统借此在一切寻根之前进行“无平方因式分解”，
这正是
\cref{exo:b1:poly:9}中无重根论证的\omterm{def:b1:poly:def}{多项式}版本。
\end{example}

\begin{theorem}[代数基本定理]\label{thm:b1:poly:fta}
$\C[X]$ 中的每个非常数\omterm{def:b1:poly:def}{多项式}在
$\C$ 中有根。\index{代数基本定理}
\end{theorem}
\admitted

\begin{remark}
尽管名字如此，该定理实为一条\emph{分析学}命题：
所有已知证明都以某种形式用到了 $\R$ 的完备性，且
没有一个是纯代数的——严格的证明放在
本科三年级分册中给出，待复积分或紧性论证
具备后方可给出。本章真正证明的是
\emph{归约}：假定每个非常数\omterm{def:b1:poly:def}{多项式}都有一个根，
则下文中 $\C$ 与 $\R$ 上的完全分解
仅凭代数即可推出。
\end{remark}

\begin{corollary}[$\C$ 与 $\R$ 上的因式分解]\label{cor:b1:poly:factorization}
\begin{enumerate}
  \item 每个非零的 $P \in \C[X]$ 可分解为
        \[
        P = c\, (X - a_1)^{m_1} \cdots (X - a_r)^{m_r},
        \]
        其中 $c$ 为首项系数，$a_i$ 为两两不同的复
        根，$\sum m_i = \deg P$：\emph{计入\omterm{def:b1:poly:derivative}{重数}后，
        一个 $n$ 次\omterm{def:b1:poly:def}{多项式}恰有 $n$ 个复根}。
  \item 每个非零的 $P \in \R[X]$ 在 $\R$ 上可分解为
        \[
        P = c \prod_i (X - a_i)^{m_i} \prod_j (X^2 + p_j X +
        q_j)^{n_j},
        \]
        其中二次因式两两不同，且 $p_j^2 - 4q_j < 0$
        （无实根）。
\end{enumerate}
\end{corollary}

\begin{proof}
(1) 对次数作归纳，每次依
\cref{thm:b1:poly:factor}分出一个根；每一步次数都相符。

(2) 设 $P$ 的系数为实数。若 $z$ 为复根，
其\omterm{def:b1:poly:derivative}{重数}为 $m$，则 $\conj z$ 亦然：对 $P(z) = 0$ 取共轭，得
$P(\conj z) = \conj{P(z)} = 0$（系数都等于其自身的共轭，
即为实数），对导数亦然
（\cref{prop:b1:poly:multiplicity}）。将非实根按
共轭成对分组，每一对贡献
\[
(X - z)(X - \conj z) = X^2 - 2\Re(z)\, X + \abs z^2 ,
\]
一个判别式为负的实二次式。实根则
贡献一次因式。
\end{proof}

\begin{example}\label{ex:b1:poly:factorexample}
$X^4 + 4$ 在 $\R$ 上的分解已见于\cref{exo:b1:complex:5}，方法是通过配对
四个复根 $\pm 1 \pm \iu$：
$X^4 + 4 = (X^2 - 2X + 2)(X^2 + 2X + 2)$。两个二次式
在 $\R$ 上均不能分解（判别式均为 $-4$）。注意：一个\emph{不可约}的实
\omterm{def:b1:poly:def}{多项式}次数为 $1$ 或 $2$——这恰是
因式分解定理所说的内容。对
$X^4 + 1$ 做同样的共轭配对，其根为 $\eu^{\pm\iu\pi/4}$ 与
$\eu^{\pm3\iu\pi/4}$：每对贡献 $X^2 -
2\cos\theta\,X + 1$，于是
\[
X^4 + 1 = \bigl(X^2 - \sqrt2\,X + 1\bigr)
\bigl(X^2 + \sqrt2\,X + 1\bigr) ,
\]
这一恒等式在 $\Q$ 上的朴素因式分解尝试中不可见——
这是坚持实（此处甚至是无理）
系数的代价，也是积分 $\frac1{x^4 +
1}$ 的标准入口，在\cref{ch:b1:integration}中。
\end{example}

\begin{omfigure}
\begin{tikzpicture}[xscale=2.6, yscale=1.35]
  \draw[->, thin] (-1.15,0) -- (1.2,0) node[below] {$x$};
  \draw[->, thin] (0,-1.35) -- (0,1.4) node[left] {$y$};
  \draw[dashed, gray] (-1,1) -- (1,1);
  \draw[dashed, gray] (-1,-1) -- (1,-1);
  \draw[gray, thin] (-1,-1.1) -- (-1,1.1) (1,-1.1) -- (1,1.1);
  \draw[omThm, very thick, domain=-1:1, samples=300, smooth]
    plot (\x, {16*\x^5 - 20*\x^3 + 5*\x});
  \node[right, font=\small] at (1.02,1) {$1$};
  \node[right, font=\small] at (1.02,-1) {$-1$};
  \fill[omThm] (-1,-1) circle (0.018) (-0.809,1) circle (0.018)
    (-0.309,-1) circle (0.018) (0.309,1) circle (0.018)
    (0.809,-1) circle (0.018) (1,1) circle (0.018);
\end{tikzpicture}

{\small \omterm{pb:b1:poly:1}{切比雪夫多项式} $T_5 = 16X^5 - 20X^3 + 5X$ 在
$\intcc{-1}1$ 上：它恰好在 $-1$ 与 $1$ 之间振荡，
在六个点（图中已标出）处触及边界值。这
\emph{等振荡}使 $2^{-4}T_5$ 成为\omterm{def:b1:poly:def}{首
一}的五次\omterm{def:b1:poly:def}{多项式}，在该区间上以上确界范数衡量最小
（\cref{exo:b1:poly:10}及周末问题）。}
\end{omfigure}

\begin{remark}[多项式的常见误区]\label{rem:b1:poly:pitfalls}
\begin{enumerate}
  \item \emph{\omterm{def:b1:poly:def}{多项式}与函数。}在 $K = \Q, \R, \C$
        上，两个概念是一致的（相等的函数有相等的系数，
        由\cref{cor:b1:poly:nroots}及 $K$ 的无限性），
        但从概念上说，\omterm{def:b1:poly:def}{多项式}就是它的系数列表：在\cref{ch:b1:structures}中的
        二元域 $\Z/2\Z$ 上，$X^2 + X$ 在两点处
        都取零值，
        但它不是零\omterm{def:b1:poly:def}{多项式}。
  \item \emph{加法下的次数。}$\deg(P + Q)$ 在首项相消时
        可能小于 $\max(\deg P, \deg Q)$；
        只有次数互不相同时，才可以放心地写
        “$\deg(P + Q) = \max(\dots)$”。
  \item \emph{正确计算根的个数。}“$n$ 个根”在
        \cref{cor:b1:poly:factorization}中是指\emph{计\omterm{def:b1:poly:derivative}{重数}、
        在 $\C$ 中}：$X^2 + 1$ 没有实根，而
        $(X-1)^2$ 只有一个不同的根，但计\omterm{def:b1:poly:derivative}{重数}后有两个。
        混淆这三种计数的陈述，
        是最常见的错误证明的来源。
  \item \emph{不可约性依赖于域。}$X^2 - 2$ 在
        $\Q$ 上不可约，在 $\R$ 上可分解；$X^2 + 1$ 在
        $\R$ 上不可约，在 $\C$ 上可分解。
        只说“不可约”而不指明系数域，
        是没有意义的。
\end{enumerate}
\end{remark}

\section{系数与根}

\begin{theorem}[韦达公式]\label{thm:b1:poly:vieta}
设 $P = X^n + c_{n-1} X^{n-1} + \dots + c_0$ 为\omterm{def:b1:poly:def}{首一多项式}，其根为
$a_1, \dots, a_n \in \C$（计\omterm{def:b1:poly:derivative}{重数}）。则\index{韦达
公式}
\[
\sum_i a_i = -c_{n-1},
\qquad
\sum_{i < j} a_i a_j = c_{n-2},
\qquad \dots, \qquad
a_1 a_2 \cdots a_n = (-1)^n c_0 ,
\]
根的第 $k$ 个对称函数为 $(-1)^k c_{n-k}$。
\end{theorem}

\begin{proof}
由\cref{cor:b1:poly:factorization}，$P = (X - a_1)\cdots(X -
a_n)$（\omterm{def:b1:poly:def}{首一}，列出全部根）。乘积按分配律展开，
每一种选法对应一项：在各因式中
选取 $X$ 或根项 $-a_i$。若在
下标为 $i_1 < \dots < i_k$ 的因式中选取根，选取 $X$ 的是
其余 $n - k$ 个因式，则贡献 $(-a_{i_1})\cdots(-a_{i_k})\,X^{n-k}$。按
$X$ 的幂合并：
\[
P = \sum_{k=0}^{n} (-1)^k
\Bigl(\sum_{i_1 < \dots < i_k} a_{i_1}\cdots a_{i_k}\Bigr)
X^{n-k} ,
\]
由 $P = \sum_k c_{n-k}X^{n-k}$ 比较系数（\omterm{def:b1:poly:def}{多项式}的系数唯一，
见\cref{def:b1:poly:def}），得 $c_{n-k} = (-1)^k
\sigma_k$，即 $\sigma_k = (-1)^kc_{n-k}$，其中 $\sigma_k$
表示上面所述的第 $k$ 个对称函数。所列三种情形为
$k = 1$、$k = 2$ 及 $k = n$。
\end{proof}

\begin{example}\label{ex:b1:poly:vieta}
对二次式 $X^2 - sX + p$：根之和为 $s$，积为 $p$ ——
前面已多次用到（见\cref{exo:b1:complex:8}）。对\omterm{def:b1:poly:def}{首一}三次式
$X^3 + aX^2 + bX + c$，其根为 $\alpha, \beta, \gamma$：
\[
\alpha + \beta + \gamma = -a,
\quad
\alpha\beta + \beta\gamma + \gamma\alpha = b,
\quad
\alpha\beta\gamma = -c ,
\]
由此无需求解即可算出 $\alpha^2 + \beta^2
+ \gamma^2 = a^2 - 2b$ 这类对称量。
\end{example}

\begin{example}[不求根的根变换]\label{ex:b1:poly:transformroots}
设 $\alpha, \beta$ 为 $X^2 - 3X + 1$ 的根。以
$\alpha^2, \beta^2$ 为根的\omterm{def:b1:poly:def}{首一}二次式是什么？由韦达公式，$\alpha +
\beta = 3$，$\alpha\beta = 1$，于是
\[
\alpha^2 + \beta^2 = (\alpha+\beta)^2 - 2\alpha\beta = 7,
\qquad
\alpha^2\beta^2 = (\alpha\beta)^2 = 1 :
\]
答案为 $X^2 - 7X + 1$ ——无需计算
$\alpha = \frac{3 + \sqrt5}2$。（验证：$\alpha^2 = \frac{7 +
3\sqrt5}2$，且确有 $\alpha^2 + \beta^2 = 7$。）同一
策略也适用于倒数（$X^2 - \frac ba X + \frac
ca$ 型变换）、平移以及任意对称数据：韦达公式把关于
\emph{未知}根的问题化为关于
\emph{已知}系数的代数运算。当根为特征值时（见
\cref{ch:b1:det}），它将经常发挥作用。
\end{example}

\begin{example}[回文方程]\label{ex:b1:poly:palindrome}
解方程 $X^4 + X^3 - 4X^2 + X + 1 = 0$。其系数从两端读
都相同，因此 $0$ 不是根，且除以
$X^2$ 不会丢失解：
\[
X^2 + X - 4 + \frac1X + \frac1{X^2} = 0 .
\]
令 $y = X + \frac1X$，则 $X^2 + \frac1{X^2} = y^2 - 2$，
方程化为
\[
y^2 + y - 6 = 0 \iff (y + 3)(y - 2) = 0 .
\]
由 $X^2 - yX + 1 = 0$ 逐个求出对应的根：当 $y = 2$ 时，$X^2 -
2X + 1 = (X - 1)^2$，得二重根 $1$；当 $y = -3$ 时，
$X^2 + 3X + 1 = 0$ 得 $X = \frac{-3 \pm \sqrt5}2$。四次式
恰有四个根（计\omterm{def:b1:poly:derivative}{重数}），符合
\cref{cor:b1:poly:factorization}的要求，通过解两个二次方程得到。该
方法适用于每一个\emph{回文}\omterm{def:b1:poly:def}{多项式}：
其根成对出现为倒数对 $\{x, 1/x\}$
（将 $X$ 换为 $1/X$ 并去分母），而 $y = X +
\frac1X$ 恰是在这类数对上保持不变的量，
从而使次数减半。
\end{example}

\begin{theorem}[拉格朗日插值]\label{thm:b1:poly:lagrange}
设 $x_0, \dots, x_n$ 为 $K$ 中互不相同的点，$y_0, \dots, y_n
\in K$。恰有一个 $P \in K[X]$，其次数 $\leq n$，满足
$P(x_i) = y_i$ 对所有 $i$ 成立，即\index{拉格朗日插值}
\[
P = \sum_{i=0}^{n} y_i\, L_i,
\qquad
L_i = \prod_{j \neq i} \frac{X - x_j}{x_i - x_j} .
\]
\end{theorem}

\begin{proof}
每个 $L_i$ 的次数为 $n$，且满足 $L_i(x_i) = 1$，$L_i(x_j) =
0$，$j \neq i$ 时（每个因式在对应的 $x_j$ 处为零）。
于是所示的 $P$ 次数 $\leq n$，且满足插值条件。唯一性：
两个次数 $\leq n$ 的插值\omterm{def:b1:poly:def}{多项式}在 $n+1$
个点 $x_i$ 处取值相同，由\cref{cor:b1:poly:nroots}可知二者相等。
\end{proof}

\begin{remark}[插曲：多项式也是向量]\label{rem:b1:poly:interlude}
一种观点上的转变，将在\cref{ch:b1:vspaces}中正式确立：
次数 $\leq n$ 的\omterm{def:b1:poly:def}{多项式}构成一个空间，其中
加法与数乘的行为恰如坐标——
\omterm{def:b1:poly:def}{多项式}\emph{就是}其 $n + 1$ 个系数构成的列表。本章中有三条结论
实则都是线性代数的内容。
\omterm{thm:b1:poly:lagrange}{拉格朗日插值}（见\cref{thm:b1:poly:lagrange}）说明
求值数据 $(P(x_0), \dots, P(x_n))$ 唯一地确定 $P$，
即在 $n + 1$ 个点处的求值是线性\omterm{def:b1:logic:inj}{双射}，
而 $L_i$ 即与之相适应的基。展开式 $R =
\sum c_k(X - a)^k$ 出现于
\cref{prop:b1:poly:multiplicity}的证明中指出，$(X - a)$ 的各次幂
构成另一套坐标系，其中 $c_k =
R^{(k)}(a)/k!$ 为坐标。又
\cref{cor:b1:poly:nroots}——根的个数多于次数会迫使
零\omterm{def:b1:poly:def}{多项式}——是一切唯一性的核心：它将变为
“在 $n + 1$ 维空间上为\omterm{def:b1:logic:inj}{单射}的线性\omterm{def:b1:logic:map}{映射}”，
并在\cref{ch:b1:findim}中正式讨论。待那些章节到来时，
空间 $K_n[X]$ 将是它们最喜爱的例子，宜
在到达那里之前便已熟练掌握它。
\end{remark}

\begin{remark}[本章的应用]\label{rem:b1:poly:whereused}
在 $\R$ 与 $\C$ 上的因式分解
（见\cref{cor:b1:poly:factorization}）是部分分式的核心，
部分分式法见\cref{ch:b1:fractions}，因而也是
在\cref{ch:b1:integration}中一大类积分的依据。将\omterm{def:b1:poly:def}{多项式}按
$(X - a)$ 的幂展开，此展开在
\cref{prop:b1:poly:multiplicity}的证明中出现，是
泰勒公式（见\cref{ch:b1:taylor}）的代数影子。特征\omterm{def:b1:poly:def}{多项式}
在微分方程中已经出现过（见\cref{ch:b1:diffeq}），
并在\cref{ch:b1:det}中重新用于矩阵。\omterm{thm:b1:poly:lagrange}{拉格朗日
插值}是数值分析中第一个存在唯一性定理，
而\omterm{pb:b1:poly:1}{切比雪夫多项式}（见\cref{exo:b1:poly:10}）
——其最优性由下文的周末问题证明——
告诉该学科应在\emph{何处}插值。
最后，$K[X]$ 的全部算术（取自
\cref{ch:b1:arith}），为
本科二年级分册中 $K[X]$ 的理想与商环的研究提供基础。
\end{remark}

\section{练习}

\begin{exercise}[$\star$]\label{exo:b1:poly:1}
作欧几里得除法：$X^5 - 1$ 除以 $X^2 + X + 1$；再
$2X^4 + X^3 - X + 3$ 除以 $X^2 - 2$。
\end{exercise}

\begin{exercise}[$\star$]\label{exo:b1:poly:2}
对哪些 $n \in \N$，$X^2 + X + 1$ \omterm{def:b1:arith:divides}{整除} $X^{2n} + X^n + 1$？
\emph{提示：$X^2 + X + 1$ 的根是 $j$ 与 $j^2$，其中 $j =
\eu^{2\iu\pi/3}$；按 $n$ 模 $3$ 分情形讨论。}
\end{exercise}

\begin{exercise}[$\star$]\label{exo:b1:poly:3}
求实数 $a, b$，使 $(X-1)^2$ \omterm{def:b1:arith:divides}{整除} $P = X^4 + aX^3 +
bX^2 + 1$，并在这些取值下将 $P$ 在 $\R$ 上因式分解。
\end{exercise}

\begin{exercise}[$\star$]\label{exo:b1:poly:4}
分别在 $\C$ 与 $\R$ 上因式分解：$X^3 - 1$；$\;X^4 + X^2 + 1$；
$\;X^6 - 1$。
\end{exercise}

\begin{exercise}[$\star\star$]\label{exo:b1:poly:5}
设 $P = X^3 - 6X^2 + 11X - 6$。
\begin{enumerate}
  \item 求有理根\emph{（\omterm{def:b1:poly:def}{首一}整系数\omterm{def:b1:poly:def}{多项式}的既约有理根 $p/q$
        是整数，且为常数项的因数，
——请证明之）}，并将 $P$ 因式分解。
  \item 不解方程，求根的平方和与倒数和，
        利用韦达定理计算，并在
        因式分解上检验结果。
\end{enumerate}
\end{exercise}

\begin{exercise}[$\star\star$]\label{exo:b1:poly:6}
用\omterm{met:b1:arith:euclid}{欧几里得算法}计算 $\gcd(X^4 - 1,\; X^3 - X^2 + X - 1)$，
并将其写成两\omterm{def:b1:poly:def}{多项式}的组合
$AU + BV$ 的形式。
\end{exercise}

\begin{exercise}[$\star\star$]\label{exo:b1:poly:7}
设 $P \in \R[X]$，使 $P(x) \geq 0$ 对所有 $x \in \R$ 成立。证明：
$P$ 是两个实\omterm{def:b1:poly:def}{多项式}的平方和：$P = A^2 + B^2$。
\emph{提示：在实因式分解中，实根的\omterm{def:b1:poly:derivative}{重数}为偶数；
将二次因式写作 $(X - z)(X - \conj z)$ 的形式，
并将 $\abs{\,\cdot\,}^2 = (\Re)^2 + (\Im)^2$ 用于
若干 $(X - z)$ 的乘积。}
\end{exercise}

\begin{exercise}[$\star\star$]\label{exo:b1:poly:8}
求\omterm{def:b1:poly:def}{多项式} $P$，使其次数 $\leq 2$，且 $P(0) = 1$，$P(1) =
3$，$P(2) = 2$。先用拉格朗日公式求解，再由关于系数的
线性方程组求解，并验证两种解法的结果一致。
\end{exercise}

\begin{exercise}[$\star\star$]\label{exo:b1:poly:9}
证明：$P = X^{2n+1} - 1$ 恰有一个实根，并且对每个 $n \geq 1$，
\omterm{def:b1:poly:def}{多项式} $1 + X + \frac{X^2}{2!} + \dots +
\frac{X^n}{n!}$ 没有重根。\emph{（比较 $P$ 与
$P'$）}。
\end{exercise}

\begin{exercise}[$\star\star\star$]\label{exo:b1:poly:10}
（\omterm{pb:b1:poly:1}{切比雪夫多项式}）定义 $T_0 = 1$，$T_1 = X$，$T_{n+1} =
2X\,T_n - T_{n-1}$。
\begin{enumerate}
  \item 用归纳法证明：$T_n(\cos\theta) = \cos n\theta$ 对
        所有 $\theta$ 成立。
  \item 由此求出 $n$ 个根（$T_n$ 的根）及其首项系数。
  \item 证明：$\sup_{x \in \intcc{-1}{1}} \abs{T_n(x)} = 1$，
        且该上确界在 $n + 1$ 个点处取到，这些点都属于 $\intcc{-1}{1}$。
\end{enumerate}
\end{exercise}

\begin{exercise}[$\star\star\star$]\label{exo:b1:poly:11}
设 $P \in \C[X]$ 为非常数\omterm{def:b1:poly:def}{多项式}，其互异根为 $a_1, \dots,
a_r$（\omterm{def:b1:poly:derivative}{重数}分别为 $m_1, \dots, m_r$）。证明有理函数的
恒等式：
\[
\frac{P'(X)}{P(X)} = \sum_{i=1}^{r} \frac{m_i}{X - a_i},
\]
由此推出高斯--卢卡斯定理：$P'$ 的每个根都落在
$P$ 的根的凸包中\emph{（代入恒等式，取 $w$ 为 $P'$ 的根，且它不是
$P$ 的根，取共轭，并把结果读作 $w$ 是
各 $a_i$ 的加权平均）}。
\end{exercise}

\begin{exercise}[$\star\star$]\label{exo:b1:poly:12}
（\omterm{def:b1:complex:unity}{单位根}滤法）设 $n \in \N^*$，$j = \eu^{2\iu\pi/3}$。
将 $(1 + X)^n$ 分别在 $1$、$j$ 与 $j^2$ 处取值，证明
\[
\sum_{k \geq 0} \binom{n}{3k}
= \frac{2^n + 2\cos\frac{n\pi}{3}}{3} ,
\]
并对 $n = 3$ 与 $n = 6$ 检验该公式。\emph{提示：$1 +
j^m + j^{2m}$ 的值为 $3$（若 $3 \mid m$），否则为 $0$；又 $1 +
j = \eu^{\iu\pi/3}$。}
\end{exercise}

\section{周末问题：切比雪夫多项式与最平坦的多项式}

\begin{problem}\label{pb:b1:poly:1}
在所有\emph{\omterm{def:b1:poly:def}{首一}}的 $n$ 次\omterm{def:b1:poly:def}{多项式}中，哪一个在
$\intcc{-1}1$ 上最接近零？答案——切比雪夫定理，
近似理论的“出生证明”——是
$2^{1-n}T_n$，其中 $T_n$ 是
\cref{exo:b1:poly:10}中的\omterm{pb:b1:poly:1}{切比雪夫多项式}，且任何\omterm{def:b1:poly:def}{首一}竞争者的偏差
都不能小于 $2^{1-n}$。\index{切比雪夫多项式}本题
展开讨论族 $(T_n)$ 的代数性质（复合律、
显式系数、第二类族 $U_n$，以及一个
微分方程），并证明极值定理及其
等号成立的情形，并列举若干应用：最优插值
节点、$\cos 36^\circ$ 的精确值，以及一个\omterm{def:b1:arith:congruence}{同余}式
$T_p \equiv X^p \pmod p$。全文约定 $T_0 = 1$，$T_1 = X$，
$T_{n+1} = 2X\,T_n - T_{n-1}$，并自由使用
$T_n(\cos\theta) = \cos n\theta$，该式由\cref{exo:b1:poly:10}得到。

\medskip
\noindent\textbf{第一部分——族 $(T_n)$。}
\begin{enumerate}
  \item 由递推关系计算 $T_2, T_3, T_4, T_5$。（比较 $T_3$ 与
        恒等式 $\cos3\theta = 4\cos^3\theta -
        3\cos\theta$，见\cref{ex:b1:complex:cos3}。）
  \item 用归纳法证明：$\deg T_n = n$，首项
        系数为 $2^{n-1}$（$n \geq 1$），且 $T_n$ 具有
        与 $n$ 相同的奇偶性（只出现偶次幂，或只出现奇次幂）。
  \item 证明唯一性原理：$T_n$ 是\emph{唯一}满足
        $P(\cos\theta) = \cos n\theta$ 对所有
        $\theta$ 成立的\omterm{def:b1:poly:def}{多项式}。（在 $\intcc{-1}1$ 上相等的两个
        \omterm{def:b1:poly:def}{多项式}处处相等，见\cref{cor:b1:poly:nroots}）。
  \item 由此推出复合律与乘积律：
        \[
        T_m \circ T_n = T_{mn},
        \qquad
        2\,T_m T_n = T_{m+n} + T_{\abs{m-n}} .
        \]
  \item 回顾\cref{exo:b1:poly:10}中的根 $x_k =
        \cos\frac{(2k+1)\pi}{2n}$ 与等振点
        $y_k = \cos\frac{k\pi}n$ 满足 $T_n(y_k) = (-1)^k$。写出
        $T_n$ 在 $\R$ 上的完全因式分解，并
        证明 $y_k$ 相互交错：$y_n < x_{n-1} < y_{n-1}
        < \dots < x_0 < y_0$。
  \item 证明：$T_n(\cosh t) = \cosh(nt)$ 对所有 $t \in \R$ 成立
        （同样用归纳法，利用
        \cref{prop:b1:functions:hyprules}），并由此对 $x
        \geq 1$ 推出闭式
        \[
        T_n(x) = \frac{\bigl(x + \sqrt{x^2 - 1}\bigr)^n +
        \bigl(x - \sqrt{x^2 - 1}\bigr)^n}{2} ,
        \]
        从而 $T_n(x) > 1$ 对 $x > 1$ 成立：在 $\intcc{-1}1$ 之外，
        \omterm{def:b1:poly:def}{多项式}立即逸出。
\end{enumerate}

\medskip
\noindent\textbf{第二部分——系数、族 $U_n$ 与
一个微分方程。}
\begin{enumerate}[resume]
  \item 由棣莫弗公式
        （见\cref{cor:b1:complex:demoivre}），证明显式
        表达式：
        \[
        T_n(x) = \sum_{0 \leq 2j \leq n} \binom{n}{2j}\,
        x^{\,n-2j}\,(x^2 - 1)^j ,
        \]
        并对 $n = 3$ 验证该式。
  \item 计算 $T_n(1)$、$T_n(-1)$ 与 $T_n(0)$（对所有 $n$）。
  \item 定义 $U_n$（\emph{第二类}）为 $U_0 = 1$，$U_1 =
        2X$，$U_{n+1} = 2X\,U_n - U_{n-1}$。证明：
        $U_n(\cos\theta) = \frac{\sin(n+1)\theta}{\sin\theta}$
        当 $\theta \notin \pi\Z$ 时成立，并且 $T_n' = n\,U_{n-1}$
        对 $n \geq 1$ 成立。
  \item 证明：$\abs{\sin n\theta} \leq n\,\abs{\sin\theta}$
        对所有 $\theta$ 成立（用归纳法），并由此推出马尔可夫型
        估计：
        \[
        \abs{T_n'(x)} \leq n^2
        \quad\text{在 } \intcc{-1}1,
        \qquad\text{其中 } T_n'(\pm1) = (\pm1)^{n-1}\,n^2 .
        \]
  \item 证明 $y = T_n$ 满足微分方程
        \[
        (1 - x^2)\,y'' - x\,y' + n^2\,y = 0 ,
        \]
        通过对恒等式 $\sin\theta\,
        T_n'(\cos\theta) = n\sin n\theta$ 关于 $\theta$ 求导，
        即可得到该方程；并直接对 $T_2$ 验证。
\end{enumerate}

\medskip
\noindent\textbf{第三部分——切比雪夫极值定理。}
设 $\widetilde T_n = 2^{1-n}\,T_n$（由第 2 题知其为\omterm{def:b1:poly:def}{首一}），并
记 $\norm{P}_\infty = \sup_{x \in \intcc{-1}1}\abs{P(x)}$。
\begin{enumerate}[resume]
  \item 证明 $\norm{\widetilde T_n}_\infty = 2^{1-n}$，该值
        在 $n + 1$ 个点 $y_n < \dots
        < y_0$ 处以正负交替的符号取得。
  \item 假设某\omterm{def:b1:poly:def}{首一多项式} $P$（次数为 $n$）满足
        $\norm P_\infty < 2^{1-n}$，令 $D = \widetilde T_n -
        P$。证明 $\deg D \leq n - 1$，并证明 $D(y_k)$
        严格与 $(-1)^k$ 同号，其中 $k = 0, \dots, n$。
  \item 由此推出 $D$ 至少有 $n$ 个互异的实根（每个间隙中一个，
        由介值性质得到，该性质在高中阶段使用，
        证明见
        \cref{ch:b1:continuity}），并得出\emph{切比雪夫
        定理}：对每个\omterm{def:b1:poly:def}{首一多项式} $P$（次数为 $n$），有
        \[
        \norm{P}_\infty \geq 2^{1-n} .
        \]
  \item （等号情形，第一步）现设 $\norm P_\infty =
        2^{1-n}$ 恰好成立，$P$ 为 $n$ 次\omterm{def:b1:poly:def}{首一多项式}，并令 $D =
        \widetilde T_n - P$。证明 $(-1)^kD(y_k) \geq 0$ 对所有 $k$ 成立，
        并证明若 $D(y_k) = 0$，且 $y_k$ 为\emph{内部}
        点（$0 < k < n$），则 $D'(y_k) = 0$ 亦成立。
        \emph{（在内部点 $y_k$ 处，$\widetilde T_n$ 与
        $P$ 都取得绝对值为 $\norm{\cdot}
        _\infty$ 的极值；可导函数在内部极值点处导数为零，
        此结论在高中阶段使用，
        证明见\cref{ch:b1:derivative}。）}
  \item （等号情形，结论）统计 $D$ 的根（计入\omterm{def:b1:poly:derivative}{重数}），以证明 $D = 0$，
        从而使范数最小的\omterm{def:b1:poly:def}{多项式}
        是\emph{唯一}的，即 $P = \widetilde T_n$。
  \item 推广到任意线段 $\intcc ab$：证明
        \omterm{def:b1:poly:def}{首一} $n$ 次\omterm{def:b1:poly:def}{多项式}的最小上确界范数
        在 $\intcc ab$ 上为 $2\bigl(\frac{b-a}4\bigr)^n$，该值由
        伸缩后的\omterm{pb:b1:poly:1}{切比雪夫多项式}取得。\emph{（作代换 $x =
        \frac{a+b}2 + \frac{b-a}2\,t$，并追踪
        首项系数。）}
\end{enumerate}

\medskip
\noindent\textbf{第四部分——应用。}
\begin{enumerate}[resume]
  \item 手算 $n = 3$ 的情形：求
        $\widetilde T_3 = X^3 - \frac34X$ 在 $\intcc{-1}1$ 上的极值点，
        验证四重等波动，其值为
        $\frac14$，并断言没有\omterm{def:b1:poly:def}{首一}三次\omterm{def:b1:poly:def}{多项式}能做得更好。
  \item （最优插值节点）对 $n + 1$ 个节点 $x_0,
        \dots, x_n \in \intcc{-1}1$，插值误差由
        $\omega(X) = \prod_i (X - x_i)$ 控制（正如
        \cref{ch:b1:taylor} 将定量刻画）。证明
        使 $\norm\omega_\infty$ 最小的选取恰为 $n + 1$ 个
        $T_{n+1}$ 的根，且 $\norm\omega_\infty = 2^{-n}$：
        切比雪夫节点正是插值的最佳位置。
  \item 利用 $T_5$，证明 $c = \cos 36^\circ$ 满足
        $16c^5 - 20c^3 + 5c + 1 = 0$，将该\omterm{def:b1:poly:def}{多项式}分解为
        $(x + 1)(4x^2 - 2x - 1)^2$，并得出
        \[
        \cos 36^\circ = \frac{1 + \sqrt5}4 .
        \]
        检验此结果与 $\cos 72^\circ =
        \frac{\sqrt5 - 1}4$ 的一致性（见\cref{exo:b1:complex:8}）。
  \item 利用第 6 题的闭式估计 $T_{10}(1.1)$，
        （两位有效数字即可），并说明：
        一个以 $1$ 为界（在 $\intcc{-1}1$ 上）的\omterm{def:b1:poly:def}{多项式}，
        其值已可超过 $40$，此时 $x = 1.1$。（$T_n$ 在此类\omterm{def:b1:poly:def}{多项式}中增长
        \emph{最快}，
        这是该族的另一条极值性质，超出本题范围。）
  \item 证明\omterm{def:b1:arith:congruence}{同余}式：对每个奇\omterm{def:b1:arith:prime}{质数} $p$，
        $T_p - X^p$ 的所有系数均被 $p$ \omterm{def:b1:arith:divides}{整除}。
        \emph{（利用第 7 题，以及 $p \mid \binom p{2j}$（$0 <
        2j < p$，该事实的证明见
        \cref{thm:b1:arith:fermat}）}在 $T_3$ 与 $T_5$ 上验证此结论。
\end{enumerate}

\medskip
\noindent\textbf{第五部分——综合。}
\begin{enumerate}[resume]
  \item 求出在 $\intcc01$ 上使上确界范数最小的\omterm{def:b1:poly:def}{首一}二次\omterm{def:b1:poly:def}{多项式}，
        并求其偏差（即第 17 题中
        $n = 2$ 的情形。）
  \item 本题究竟在何处用到了：（i）\omterm{def:b1:poly:def}{多项式}的刚性
        （\cref{cor:b1:poly:nroots}）；（ii）
        三角学（\cref{ch:b1:complex}
        与\cref{ch:b1:functions}）；（iii）\omterm{def:b1:counting:objects}{二项式系数}的算术，
        它来自\cref{ch:b1:arith}？各用一句话作答。
  \item 综述（用一小段话）：定理断言，
        最“平坦”的\omterm{def:b1:poly:def}{首一多项式}恰是
        \emph{等波动}的那一个，证明则将最优性转化为
        根的计数。请评述此机制，并说明
        代换 $x = \cos\theta$ 作为代数与三角之间桥梁的作用，
        并指出本题在哪两处用到了分析学事实
        （介值性质与内部极值），这些事实将在
        后续章节中予以证明。
\end{enumerate}
\end{problem}
